% chapters/12/object-model.tikz
\section{GameObject dan Komponen}
\label{sec:gameobject}

Benda pada Unity tidak disusun dengan mewarisi kelas, melainkan dengan
menempelkan komponen. Bagian ini menjelaskan gagasan penyusunan itu, sebab
gagasan itulah yang membedakannya dari blok Python.

\subsection{Benda Kosong yang Disusun dari Bagian}

Sebuah \istilah{GameObject} pada dasarnya adalah wadah kosong. Ia tidak
memiliki gambar, tidak memiliki bentuk, dan tidak melakukan apa pun. Yang
membuatnya berguna adalah \textbf{komponen} yang menempel padanya.

Sebuah benda yang tampak di layar tersusun atas sekurang-kurangnya tiga
komponen.

\begin{itemize}
  \item \istilah{Transform}, yang memuat letak, putaran, dan ukuran benda.
        Komponen ini wajib ada pada setiap \istilah{GameObject}.
  \item \istilah{Mesh Renderer} atau \istilah{Sprite Renderer}, yang menentukan
        bentuk atau gambar yang ditampilkan.
  \item \istilah{Collider}, yang menentukan bentuk batas benda untuk keperluan
        tabrakan.
\end{itemize}

Selain ketiganya, komponen lain dapat ditambahkan sesuai keperluan, misalnya
\istilah{Rigidbody} untuk benda yang mengikuti fisika, \istilah{AudioSource}
untuk benda yang bersuara, atau \istilah{script} C\# yang Anda tulis sendiri.

\begin{catatabox}
Komponen yang Anda tulis sendiri pun berupa \istilah{script} C\#.
Skrip itu ditempelkan pada sebuah \istilah{GameObject}, dan sejak saat itu ia
menjadi bagian dari benda tersebut. Karena itu, satu skrip yang sama dapat
ditempelkan pada banyak benda, dan masing-masing benda menjalankannya sendiri
dengan nilainya sendiri. Inilah padanan dari kelas dan objek pada blok Python,
tetapi dengan cara yang berbeda.
\end{catatabox}

\subsection{Perbandingan dengan Cara Pygame}

Perhatikan perbedaan keduanya.

\begin{itemize}
  \item \textbf{Pada Pygame}, sebuah benda adalah objek dari kelas yang Anda
        tulis. Seluruh sifat dan perilakunya ditentukan oleh kelas itu, dan
        benda yang berbeda perilakunya menuntut kelas yang berbeda.
  \item \textbf{Pada Unity}, sebuah benda adalah wadah yang diisi komponen.
        Perilakunya ditentukan oleh komponen apa saja yang menempel padanya,
        dan benda yang berbeda perilakunya cukup diberi komponen yang berbeda.
\end{itemize}

Perbedaan itu membawa akibat praktis yang penting: menambahkan kemampuan pada
sebuah benda tidak menuntut perubahan kode, melainkan cukup dengan menempelkan
komponen. Sebuah benda dapat diubah menjadi benda yang mengikuti fisika hanya
dengan menambahkan \istilah{Rigidbody}, tanpa menyentuh satu baris pun kode
yang sudah ada.

\subsection{Pengaturan Komponen pada \istilah{Inspector}}

Setiap komponen memiliki pengaturan, dan seluruhnya ditampilkan pada
\istilah{Inspector}. Pengaturan itu dapat diubah dengan mengetikkan angkanya
langsung, atau dengan menggeser penunjuk sambil menekan tetikus.

Perhatikan dua hal pada saat mengubah pengaturan.

\begin{enumerate}
  \item \textbf{Nilai yang diubah ketika scene sedang berjalan tidak
        tersimpan.} Unity sengaja memisahkan nilai yang diubah ketika berjalan
        dari nilai yang tersimpan, sehingga percobaan tidak merusak susunan
        Anda. Nilai itu akan dikembalikan ketika Anda berhenti menjalankan
        scene.
  \item \textbf{Pengaturan komponen yang tampak pada \istilah{Inspector}
        berasal dari kode Anda.} Perubahan pada nilai publik di dalam skrip
        akan tampak pada panel itu, dan nilai yang Anda ubah pada panel itu
        akan menimpa nilai bawaan pada skrip.
\end{enumerate}

\begin{tipsbox}
Kebiasaan yang berguna: berikan nilai bawaan yang masuk akal pada skrip Anda,
sehingga benda yang baru saja ditempeli skrip itu langsung berperilaku wajar
tanpa perlu diatur. Cara itu menghemat waktu, terutama ketika Anda menempelkan
skrip yang sama pada banyak benda.
\end{tipsbox}
